% ══════════════════════════════════════════════════════════════════
%  Section 4 — System Design
% ══════════════════════════════════════════════════════════════════

\section{System Design}\label{sec:system}

This section describes the four principal subsystems of NEOS---the
command interface, the field dynamics engine, the multi-field
orchestrator, and the interface layer---followed by a walkthrough of the
interactive shell.

% ─────────────────────────────────────────────────────────────────
\subsection{Command Interface}\label{sec:command-interface}

Every interaction with NEOS follows a single syntactic template:

\medskip
\begin{center}
\texttt{nf <command> [subcommand] [arguments] [-{}-flags]}
\end{center}
\medskip

\noindent
Patterns are referenced with the \texttt{@name} sigil; fields with
\texttt{\$field}.  The design is deliberately Unix-flavoured:
\emph{imperative verbs acting on semantic objects}.  Commands are
organised into six categories:

\begin{enumerate}[nosep]
  \item \textbf{Operations} --- create, inject, remove, and transform
        patterns.
  \item \textbf{Dynamics} --- advance the field through cycles, control
        decay and resonance parameters.
  \item \textbf{Measurement} --- query activation, coherence, energy,
        and attractor membership.
  \item \textbf{Visualisation} --- render field state as landscapes,
        graphs, and heatmaps.
  \item \textbf{Versioning} --- commit, branch, and diff field states.
  \item \textbf{Autonomy} --- define tasks that run to a convergence
        criterion without user intervention.
\end{enumerate}

% ── TikZ: Command pipeline ──────────────────────────────────────
\begin{figure}[htbp]
\centering
\begin{tikzpicture}[
  stage/.style={
    draw, rounded corners=6pt, fill=neosblue!6, draw=neosblue!35,
    minimum width=2.4cm, minimum height=0.9cm,
    font=\small\sffamily, align=center, thick
  },
  arr/.style={-{Stealth[length=5pt]}, thick, neosblue!40},
]

\node[stage] (inp)    at (0,0)    {User Input};
\node[stage] (parse)  at (3.2,0)  {Parser};
\node[stage] (ast)    at (6.4,0)  {Canonical\\[-2pt]AST};
\node[stage] (eng)    at (9.6,0)  {Field\\[-2pt]Engine};
\node[stage] (out)    at (12.8,0) {State Update\\[-2pt]\& Output};

\draw[arr] (inp)   -- (parse);
\draw[arr] (parse) -- (ast);
\draw[arr] (ast)   -- (eng);
\draw[arr] (eng)   -- (out);

\end{tikzpicture}
\caption{Command pipeline.  Every interaction, regardless of input
  modality, is parsed into the same Canonical AST before reaching the
  field engine.}
\label{fig:command-pipeline}
\end{figure}

% ─────────────────────────────────────────────────────────────────
\subsection{Field Dynamics}\label{sec:field-dynamics}

Each invocation of \texttt{nf cycle} advances the field through six
sequential phases:

\begin{enumerate}[nosep]
  \item \textbf{Decay.}\quad $A \leftarrow A \times (1 - \lambda)$.
        Every pattern loses a fixed fraction of its activation, ensuring
        that unreinforced ideas fade.
  \item \textbf{Resonate.}\quad Pairwise resonance is computed across
        semantic, logical, and contextual dimensions.
  \item \textbf{Amplify.}\quad
        $A \leftarrow A + \alpha \sum (R \cdot A)$.
        Patterns that resonate strongly with their neighbours gain
        activation energy.
  \item \textbf{Threshold.}\quad $A < \tau \;\Rightarrow\;
        \text{mark dormant}$.  Patterns whose activation has fallen
        below a critical value are flagged as dormant and excluded from
        further resonance computation.
  \item \textbf{Coherence.}\quad
        $C = \mu_R\, / \,(1 + \sigma^{2}_R)$.
        A scalar summary of how tightly the surviving patterns agree.
  \item \textbf{Attractor.}\quad Emergence is declared when three
        conditions hold simultaneously:
        coherence $> 0.6$, energy $> 70\%$ of maximum, and
        perturbation-stability under small noise.
\end{enumerate}

% ── TikZ: Six-phase dynamics ring ────────────────────────────────
\begin{figure}[htbp]
\centering
\begin{tikzpicture}[
  phase/.style={
    draw, circle, minimum size=2.0cm, thick,
    font=\scriptsize\sffamily, align=center, inner sep=2pt
  },
  arr/.style={-{Stealth[length=5pt]}, thick, gray!60},
]

% Place six nodes in a hexagonal ring (starting from top, clockwise)
\def\R{2.8}
\node[phase, fill=neosred!8,    draw=neosred!40]
  (p1) at (90:\R)  {Decay\\$A(1{-}\lambda)$};
\node[phase, fill=neosorange!10, draw=neosorange!50]
  (p2) at (30:\R)  {Resonate\\$R(x,y)$};
\node[phase, fill=neosgreen!8,  draw=neosgreen!40]
  (p3) at (-30:\R) {Amplify\\$A{+}\alpha\Sigma$};
\node[phase, fill=neosblue!8,   draw=neosblue!35]
  (p4) at (-90:\R) {Threshold\\$A<\tau$};
\node[phase, fill=neosviolet!8, draw=neosviolet!40]
  (p5) at (210:\R) {Coherence\\$\mu_R/(1{+}\sigma^2)$};
\node[phase, fill=neosorange!8, draw=neosorange!45]
  (p6) at (150:\R) {Attractor\\$C{>}0.6$};

% Clockwise arrows
\draw[arr] (p1.east)  to[bend left=15] (p2.north);
\draw[arr] (p2.south) to[bend left=15] (p3.north east);
\draw[arr] (p3.south west) to[bend left=15] (p4.east);
\draw[arr] (p4.west)  to[bend left=15] (p5.south);
\draw[arr] (p5.north) to[bend left=15] (p6.south west);
\draw[arr] (p6.north east) to[bend left=15] (p1.west);

\end{tikzpicture}
\caption{The six-phase dynamics ring executed during each
  \texttt{nf cycle} invocation.  Arrows indicate the fixed evaluation
  order.}
\label{fig:dynamics-ring}
\end{figure}

% ─────────────────────────────────────────────────────────────────
\subsection{Multi-Field Orchestration}\label{sec:multi-field}

A single field captures one perspective.  Real problems demand several.
NEOS supports multiple coupled fields---for instance, a
\emph{Technical} field, a \emph{Business} field, and a \emph{User}
field---each evolving under its own dynamics but linked through a
coupling matrix~$\Gamma$.

Stability requires that the spectral radius of~$\Gamma$ remain bounded:
\begin{equation}\label{eq:coupling-stability}
  \rho(\Gamma) \;<\; \frac{\lambda_{\min}}{\alpha_{\max}}\,,
\end{equation}
where $\lambda_{\min}$ is the smallest decay rate across fields and
$\alpha_{\max}$ is the largest amplification factor.  When
\eqref{eq:coupling-stability} is satisfied, cross-field influence is
strong enough to exchange information yet too weak to destabilise any
individual field.

Two orchestration modes are provided:

\begin{itemize}[nosep]
  \item \textbf{Pipeline} --- fields execute sequentially, each seeding
        the next with its collapsed output.
  \item \textbf{Parallel} --- fields execute simultaneously, exchanging
        activation through~$\Gamma$ at every cycle.
\end{itemize}

% ── TikZ: Multi-field orchestration ──────────────────────────────
\begin{figure}[htbp]
\centering
\begin{tikzpicture}[
  field/.style={
    draw, rounded corners=6pt, minimum width=2.6cm,
    minimum height=1.0cm, thick, font=\small\sffamily,
    align=center
  },
  synth/.style={
    draw, rounded corners=6pt, minimum width=3.0cm,
    minimum height=1.0cm, thick, font=\small\sffamily\bfseries,
    fill=neosgreen!8, draw=neosgreen!40, align=center
  },
  arr/.style={-{Stealth[length=5pt]}, thick, gray!60},
]

\node[field, fill=neosblue!8,   draw=neosblue!35]   (tech) at (-3.4, 2) {Technical\\Field};
\node[field, fill=neosorange!8, draw=neosorange!45]  (biz)  at ( 0.0, 2) {Business\\Field};
\node[field, fill=neosviolet!8, draw=neosviolet!40]  (usr)  at ( 3.4, 2) {User\\Field};

\node[synth] (syn) at (0, -0.4) {Synthesis};

\draw[arr] (tech.south) -- (syn.north west);
\draw[arr] (biz.south)  -- (syn.north);
\draw[arr] (usr.south)  -- (syn.north east);

% Coupling arcs between fields
\draw[dashed, gray!50, thick, <->]
  (tech.east) -- (biz.west)
  node[midway, above, font=\scriptsize\sffamily] {$\Gamma$};
\draw[dashed, gray!50, thick, <->]
  (biz.east) -- (usr.west)
  node[midway, above, font=\scriptsize\sffamily] {$\Gamma$};

\end{tikzpicture}
\caption{Multi-field orchestration.  Three perspective fields are
  coupled through~$\Gamma$ and converge into a single synthesis node.}
\label{fig:multi-field}
\end{figure}

% ─────────────────────────────────────────────────────────────────
\subsection{Interface Modes}\label{sec:interfaces}

NEOS exposes three input modalities, all of which compile to the same
Canonical AST:

\begin{enumerate}[nosep]
  \item \textbf{Semantic interface} --- natural-language statements
        parsed into field operations
        (e.g., \emph{``inject security with strength 0.9''}).
  \item \textbf{Algebraic interface} --- formal mathematical notation:
        $\iota(\text{``security''},\, 0.9)$.
  \item \textbf{Shell interface} --- the command line:
        \texttt{nf inject "security" 0.9}.
\end{enumerate}

\noindent
The three modalities are isomorphic: every expression in one can be
mechanically translated to the other two.  This triple interface ensures
that NEOS is equally accessible to domain experts (semantic), theorists
(algebraic), and engineers (shell).

% ─────────────────────────────────────────────────────────────────
\subsection{The NEOS Shell}\label{sec:shell}

The NEOS shell is a read-eval-print loop (REPL) for cognitive
computing---what \texttt{bash} is to Unix, the NEOS shell is to the
neural field.  A typical session proceeds through three phases:

\begin{description}[nosep, leftmargin=2em, labelwidth=1.8em, style=nextline]
  \item[\textnormal{\textsc{Create}}]
    Inject patterns into the field
    (\texttt{nf inject}).
  \item[\textnormal{\textsc{Compile}}]
    Collapse the superposition into a definite output
    (\texttt{nf collapse}).
  \item[\textnormal{\textsc{Run}}]
    Execute an autonomous task against a convergence criterion
    (\texttt{nf task}).
\end{description}

\noindent
Debugging a thought proceeds exactly as debugging a program: set
checkpoints (\texttt{nf checkpoint}), enter step mode
(\texttt{nf cycle 1 -{}-trace}), and inspect the cycle trace to
understand why a pattern gained or lost activation.

% ─────────────────────────────────────────────────────────────────
\subsection{Getting Started: A Walkthrough}\label{sec:tutorial}

Using NEOS requires no installation.  The operator copies the kernel
prompt (\texttt{prompts/nfos-kernel.md}) and pastes it as the system
prompt in any LLM interface that supports one.  The LLM becomes the
NEOS runtime.  The following walkthrough demonstrates the complete
analysis cycle using a startup-idea evaluation as the domain.

\paragraph*{Step~1: Create a Session.}

\begin{terminalbox}[Startup Analysis --- Session Creation]
\begin{lstlisting}[style=neosshell]
nf session new "Startup Idea Analysis"
[SESSION] Created: Startup Idea Analysis
  Field: main (default)
  Parameters: l=0.05  a=0.30  t=0.40  s=0.50
  Mode: step (interactive)
\end{lstlisting}
\end{terminalbox}

\noindent
A session initialises a default field with standard
parameters.  The operator can tune parameters at any time with
\texttt{nf tune}.

\paragraph*{Step~2: Inject Patterns.}

\begin{terminalbox}[Startup Analysis --- Pattern Injection]
\begin{lstlisting}[style=neosshell]
nf inject "strong_founding_team" 0.9
[INJECT] strong_founding_team (s: 0.90)

nf inject "crowded_market" 0.7
[INJECT] crowded_market (s: 0.70)
  [RESONANCE] <-> @strong_founding_team: R = 0.42 MODERATE

nf inject "novel_technology" 0.85
[INJECT] novel_technology (s: 0.85)

nf inject "unclear_revenue_model" 0.6
[INJECT] unclear_revenue_model (s: 0.60)

nf inject "early_traction_with_users" 0.8
[INJECT] early_traction_with_users (s: 0.80)
  [RESONANCE] <-> @strong_founding_team: R = 0.78 STRONG
\end{lstlisting}
\end{terminalbox}

\noindent
Each pattern is an observation with an activation strength
(0.0--1.0).  Higher strength indicates greater initial confidence.
Resonance is computed automatically at injection time.

\paragraph*{Step~3: Run Dynamics.}

\begin{terminalbox}[Startup Analysis --- Five Cycles]
\begin{lstlisting}[style=neosshell]
nf cycle 5 --trace
[CYCLE 1] C = 0.38
  @strong_founding_team: 0.90 -> 0.94  (resonance boost)
  @crowded_market:       0.70 -> 0.67  (decay, low resonance)
[CYCLE 2] C = 0.52
  Cluster forming: {team, traction, technology}
[CYCLE 3] C = 0.64
  [ATTRACTOR-EMERGED] "strong_fundamentals"
    Core: {@strong_founding_team, @early_traction,
           @novel_technology}
    Coherence: 0.64
[CYCLE 4] C = 0.71
[CYCLE 5] C = 0.76
  @unclear_revenue_model: 0.52 (weakening, low resonance)
  @crowded_market: 0.48 (near threshold)
\end{lstlisting}
\end{terminalbox}

\noindent
Each cycle applies the master equation: decay erodes unreinforced
patterns, resonance amplifies related ones.  The positive factors
(team, technology, traction) cluster into an attractor; risks
(crowded market, unclear revenue) weaken due to low mutual resonance.

\clearpage
\paragraph*{Step~4: Collapse to Insight.}

\begin{terminalbox}[Startup Analysis --- Collapse]
\begin{lstlisting}[style=neosshell]
nf collapse --strategy attractor
[COLLAPSE] Strategy: attractor
  Dominant: strong_fundamentals (C=0.76)

OUTPUT:
  The field identifies strong technical and team fundamentals
  with early market validation.  Primary risks are an unclear
  revenue model and a crowded market---both weakening under
  dynamics, suggesting they are less structurally connected to
  the core thesis than the positive signals.

  Confidence: HIGH (0.76)
\end{lstlisting}
\end{terminalbox}

\noindent
Collapse resolves the field's superposition into structured output.
The operator can also query the field in natural language
(\texttt{nf ask "What is the biggest risk?"}), save the state
(\texttt{nf commit}), branch to explore alternatives
(\texttt{nf branch create}), or switch to autonomous mode
(\texttt{nf mode auto}) for hands-free analysis.

\medskip
\noindent
This four-step cycle---inject, run dynamics, observe, collapse---is the
fundamental workflow of NEOS.  The empirical
evidence~(\S\ref{sec:evidence}) demonstrates the same cycle operating
at scale: 69~patterns over 52~cycles in software quality, and
10~patterns over 12~cycles in network forensics.
